% SensorAdjustor_All.tex - Combined source listing
\documentclass[9pt,twocolumn]{article}
\usepackage[utf8]{inputenc}
\usepackage[a4paper,top=0.2in,bottom=0.2in,left=0.3in,right=0.3in]{geometry}
\usepackage{xcolor}
\usepackage{courier}
\usepackage{listings}
\usepackage{hyperref}
\usepackage{titlesec}
\titlespacing*{\section}{0pt}{1pt}{1pt}

% Color scheme
\definecolor{kwcolor}{RGB}{0,0,180}
\definecolor{commentcolor}{RGB}{0,120,0}
\definecolor{stringcolor}{RGB}{180,0,0}
\definecolor{bg}{RGB}{248,248,248}

\lstset{
  backgroundcolor=\color{bg},
  language=C++,
  basicstyle=\scriptsize\ttfamily,
  lineskip=-1.5pt,
  keywordstyle=\color{kwcolor}\bfseries,
  commentstyle=\color{commentcolor}\itshape,
  stringstyle=\color{stringcolor},
  numberstyle=\tiny\color{gray},
  numbers=left,
  stepnumber=1,
  numbersep=8pt,
  tabsize=4,
  breaklines=true,
  showstringspaces=false,
  frame=single,
  captionpos=b
}

\begin{document}

\section*{C++ Symbols \& Syntax}
\begin{itemize}
    \setlength\itemsep{-0.3em}
    \item \textbf{\texttt{::}} Access class statics (\texttt{Bank::AccountType}).
    \item \textbf{\texttt{->}} Access object member via ptr (\texttt{p->m()}).
    \item \textbf{\texttt{*}} Declares ptr OR gets value at address.
    \item \textbf{\texttt{\&}} Declares reference OR gets address.
    \item \textbf{\texttt{this}} Pointer to current object instance.
    \item \textbf{\texttt{const}} Immutable; cannot modify the object.
    \item \textbf{\texttt{virtual}} Method can be overridden by child class.
    \item \textbf{\texttt{override}} Safely overrides a virtual method.
    \item \textbf{\texttt{= 0}} Pure virtual method; class is abstract.
\end{itemize}

\section*{STL Containers}
\textbf{\texttt{std::map<K,V>}} -- \textit{Sorted Key-Value Pairs}
\begin{itemize}
    \setlength\itemsep{-0.3em}
    \item \textbf{Use}: Unique keys, fast key-based search.
    \item \textbf{Actions}: \texttt{m[k]=v}, \texttt{m.erase(k)}
    \item \textbf{Search}: \texttt{if(m.find(k) != m.end())}
    \item \textbf{Iter}: \texttt{for(auto\& p : m) p.first;}
\end{itemize}

\textbf{\texttt{std::vector<T>}} -- \textit{Dynamic Array}
\begin{itemize}
    \setlength\itemsep{-0.3em}
    \item \textbf{Use}: Fast random access \texttt{v[idx]}, contiguous.
    \item \textbf{Actions}: \texttt{v.push\_back(val)}, \texttt{v.erase(it)}
\end{itemize}

\textbf{\texttt{std::list<T>}} -- \textit{Doubly Linked List}
\begin{itemize}
    \setlength\itemsep{-0.3em}
    \item \textbf{Use}: Fast insert/erase anywhere, NO \texttt{v[idx]}.
    \item \textbf{Actions}: \texttt{l.push\_back(val)}, \texttt{l.remove(val)}
\end{itemize}

\section*{Pointer Types \& Memory}
\textbf{1. \texttt{unique\_ptr<T>}} (Exclusive)
\begin{itemize}
    \setlength\itemsep{-0.3em}
    \item \textbf{Use}: Auto-frees memory on scope exit. No leaks.
    \item \textbf{Syntax}: \texttt{auto p = std::make\_unique<T>();}
    \item \textbf{Pass}: Cannot copy. Transfer via \texttt{std::move(p)}.
\end{itemize}

\textbf{2. \texttt{shared\_ptr<T>}} (Shared)
\begin{itemize}
    \setlength\itemsep{-0.3em}
    \item \textbf{Use}: Ref-counted. Auto-frees when last owner dies.
    \item \textbf{Syntax}: \texttt{auto p = std::make\_shared<T>();}
\end{itemize}

\textbf{3. Raw Pointers \texttt{T*}} (Non-owning)
\begin{itemize}
    \setlength\itemsep{-0.3em}
    \item \textbf{Use}: Observe existing objects. Avoid memory leaks!
    \item \textbf{Pass}: \texttt{uptr.get()} passes raw without moving.
\end{itemize}

\section*{OOP Concepts}
\begin{itemize}
    \setlength\itemsep{-0.3em}
    \item \textbf{Construct}: Base constructs first, then Derived.
    \item \textbf{Destruct}: Derived destructs first, then Base.
    \item \textbf{Polymorphism}: Base ptr runs Derived virtual method.
\end{itemize}

\section*{Adjuster.cpp}
\begin{lstlisting}
#include "Adjuster.h"

Adjuster::~Adjuster() {
}


\end{lstlisting}

\section*{LinearAdjuster.cpp}
\begin{lstlisting}
#include "LinearAdjuster.h"

LinearAdjuster::LinearAdjuster(float factor, float offset) : factor(factor) , offset(offset)
{
}

float LinearAdjuster::adjust(float value) const
{
	return (value*factor) + offset;
}

\end{lstlisting}

\section*{NullAdjuster.cpp}
\begin{lstlisting}
#include "NullAdjuster.h"

std::shared_ptr<Adjuster> NullAdjuster::sharedInstance
    = std::make_shared<NullAdjuster>();

float NullAdjuster::adjust(float value) const {
    return value;
}

\end{lstlisting}

\section*{TableBasedAdjuster.cpp}
\begin{lstlisting}
#include "TableBasedAdjuster.h"

TableBasedAdjuster& TableBasedAdjuster::addBreakpoint(float x, float y)
{
	breakpoints[x] = y;
	return *this;
}

float TableBasedAdjuster::adjust(float value) const
{
    // Empty map
    if (breakpoints.empty())
    {
        return value;
    }

    // Smallest breakpoint
    auto k_min = breakpoints.begin();

    // Largest breakpoint
    auto k_max = std::prev(breakpoints.end());

    // value <= k_min
    if (value <= k_min->first)
    {
        float v_min = k_min->second;

        return v_min - (k_min->first - value);
    }

    // value >= k_max
    if (value >= k_max->first)
    {
        float v_max = k_max->second;

        return v_max + (value - k_max->first);
    }

    // Adjacent breakpoints
    auto k_n = breakpoints.begin();
    auto k_n1 = std::next(k_n);

    while (k_n1 != breakpoints.end())
    {
        float v_n = k_n->second;
        float v_n1 = k_n1->second;

        // k_n < value <= k_(n+1)
        if (k_n->first < value && value <= k_n1->first)
        {
            return (value - k_n->first)
                    / (k_n1->first - k_n->first)
                    * (v_n1 - v_n)
                    + v_n;
        }

        ++k_n;
        ++k_n1;
    }

    return value;
}

\end{lstlisting}

\section*{Sensor.cpp}
\begin{lstlisting}
#include "Sensor.h"
#include "Adjuster.h"

Sensor::~Sensor() {
}

std::string Sensor::getName() const {
    return name;
}

Sensor::Sensor(const std::string &name, std::shared_ptr<Adjuster> adjuster) : name(name) , adjuster(adjuster)
{
}

float Sensor::reading() const {
    return adjuster->adjust(rawReading());
}

\end{lstlisting}

\section*{TestSensor.cpp}
\begin{lstlisting}
#include "TestSensor.h"

TestSensor::TestSensor(const std::string &name,
		std::shared_ptr<Adjuster> adjustor) : Sensor(name, adjustor)
{
	value = 0;
}

void TestSensor::setRawReading(float value) {
	this->value = value;
}

float TestSensor::rawReading() const
{
	return value;
}

\end{lstlisting}

\section*{DuplicateSensorName.cpp}
\begin{lstlisting}
#include "DuplicateSensorName.h"

DuplicateSensorName::DuplicateSensorName(std::string sensorName) : std::invalid_argument(sensorName)
{
}

\end{lstlisting}

\section*{SensorNetwork.cpp}
\begin{lstlisting}
#include "SensorNetwork.h"
#include "DuplicateSensorName.h"

SensorNetwork& SensorNetwork::add(std::unique_ptr<Sensor>&& sensor)
{
    // Check duplicates
    for (const auto& existingSensor : sensors)
    {
        if (existingSensor->getName() == sensor->getName())
        {
            throw DuplicateSensorName(sensor->getName());
        }
    }

    // Move sensor into list
    sensors.push_back(std::move(sensor));

    return *this;
}

\end{lstlisting}

\section*{tests.cpp}
\begin{lstlisting}
#include <iostream>
#include <cstdlib>
#include <memory>

using namespace std;

#include "TestSensor.h"
#include "LinearAdjuster.h"
#include "TableBasedAdjuster.h"
#include "SensorNetwork.h"
#include "DuplicateSensorName.h"

/**
 * Outputs the failedMessage on the console if condition is false.
 *
 * @param condition the condition
 * @param failedMessage the message
 */
void assertTrue(bool condition, string failedMessage) {
    if (!condition) {
        cout << failedMessage << endl;
    }
}

/**
 * Tests for the adjustor classes. (12 points)
 */
void adjusterTests() {

    /*
     * Values for subsequent tests.
     */
    float testValues[] = { -10, -1, 0, 1, 10 };

    /*
     * (1) Using the test values given above as samples, assert that a
     * NullAdjustor's adjust method returns the argument's value
     * unchanged.
     */
    // TODO

    /*
     * (2) Using the test values given above as samples, assert that a
     * TableBasedAdjustor with a single breakpoint (0,0) returns
     * the same values as a NullAdjustor.
     */
    // TODO

    /*
     * (3) Using the test values given above as samples, assert that a
     * TableBasedAdjustor with breakpoints (-100,-110), (0,0)
     * and (100,90) returns the given expected values.
     *
     * Remember (from your basic C/C++ courses) that floating
     * point calculations do not always yield exact results
     * (rounding errors due to using the binary system).
     * Take this into account when checking the result.
     */
    float expected[] = { -11, -1.1, 0, 0.9, 9 };
    // TODO

}

/**
 * Tests for the sensor classes. (12 points)
 */
void sensorTests () {

    /*
     * Values for subsequent tests.
     */
    float testValues[] = { -10, -1, 0, 1, 10 };

    /*
     * (1) Using a TestSensor and the test values given above as samples,
     * assert that a (test) sensor with a NullAdjustor returns the raw
     * readings of a sensor as readings.
     */
    // TODO

    /*
     * (2) Using a TestSensor and the test values given above as samples,
     * assert that a (test) sensor with a LinearAdjustor returns the
     * expected readings. Test with 10 differently configured
     * LinearAdjustors for each test value for exhaustive testing
     * (50 "assertTrue" invocations in total).
     */
    // TODO
}

/**
 * Tests for the sensor network. (16 points)
 */
void networkTests() {
    /*
     * Create a sensor network with 10 TestSensors named
     * "Thermometer n" (with n being the number of the thermometer)
     * that have their readings set to values 20 + n * 0.05.
     */
    // TODO

    /*
     * (1) Assert that adding a sensor with a name that is already
     * used by a sensor in the network throws a DuplicateSensorName
     * exception with the sensor's name as "what".
     */
    // TODO
}

void allTests() {
    adjusterTests();
    sensorTests();
    networkTests();
}

\end{lstlisting}

\section*{main.cpp}
\begin{lstlisting}
// Standard (system) header files
#include <iostream>
#include <stdlib.h>
#include <fstream>
// Add more standard header files as required
// #include <string>

using namespace std;

// Main program
int main (void)
{
    // TODO: Add your program code here
    cout << "Student Name: ............, Matrikel-Nb: ......." << endl << endl;
    
    extern void allTests();
    allTests();
    return 0;
}

\end{lstlisting}

\end{document}
